\section*{ID9205: Definition of θi}
\paragraph{Metadata.}
PiID: 1032748892186 | RTEID: RTE:P35745.0592:T2.3482 | UUID: d4aa9bda-1599-5b0b-a5b4-bad8a8f2742e

\paragraph{Canonical Statement.}
and frequency ratios into canonical TZPQVT harmonic values. Falsifiability / Test Handle: Confirm by comparing the node counts in each system and verifying that ratios of measured resonant frequencies or lag times scale by approximately 1.1851851851.1851851851.185185185 when mapped between frameworks. 2.0.152 ID 301: Fibonacci Sphere Parametrization Type: Definition Formal Statement: For a target of n=153,600n=153\{,\}600n=153,600 points on the unit sphere S2Sˆ2S2, the iiith address is given by \$θi=2πi/φ\textbackslash{}theta\_i\$ = 2\textbackslash{}pi \$i/\textbackslash{}varphiθi=2πi/φ\$ and i=arccos(12(i0.5)/n)\textbackslash{}phi\_i = \textbackslash{}arccos\textbackslash{}bigl(1 - 2(i - 0.5)/n\textbackslash{}bigr)i=arccos(12(i0.5)/n), with \$φ=(1+5)/2\textbackslash{}varphi\$ = \$(1+\textbackslash{}sqrt\{5\})/2φ=(1+5)/2\$ the golden ratio. The Cartesian coordinates are \$pi=(sinicosθi,sinisinθi,cosi)p\_i\$ = \$(\textbackslash{}sin\textbackslash{}phi\_i\textbackslash{}cos\textbackslash{}theta\_i,\textbackslash{}sin\textbackslash{}phi\_i\textbackslash{}sin\textbackslash{}theta\_i,\textbackslash{}cos\textbackslash{}phi\_i)pi=(sinicosθi,sinisinθi,cosi).\$ Interpretive Meaning: This algorithm uniformly di

\paragraph{Canonical Equation.}
\[
θi=2πi/φ\theta_i
\]

\paragraph{Dictionary.}
Definition of θi — θi=2πi/φθ\_i

\paragraph{Encyclopedia.}
Definition of θi is an algebraic definition, written θi=2πi/φθ\_i. It is an algebraic definition expressing the quantity directly in terms of the variables on its right-hand side. Within the Trawin Zero Point registry it serves as one element of the framework's mathematical narrative.
